\subsection{Localización de la planitud}
\label{subsection:0.6.3-fr}

\begin{env}[6.3.1]
\label{0.6.3.1-fr}
Si $A$ es un anillo, $S$ una parte multiplicativa de $A$, $S^{-1}A$ es
un $A$-\emph{módulo plano}. En efecto, para todo $A$-módulo $N$,
$N\otimes_A S^{-1}A$ se identifica con $S^{-1}N$ (1.2.5) y se sabe (1.3.2)
que $S^{-1}N$ es un funtor exacto en $N$.

Si ahora $M$ es un $A$-módulo plano,
$S^{-1}M=M\otimes_A S^{-1}A$ es un $S^{-1}A$-módulo plano (6.2.1), por tanto
es también $A$-plano en virtud de lo que precede y de (6.2.1). En particular,
si $P$ es un $S^{-1}A$-módulo, se lo puede considerar como un $A$-módulo
isomorfo a $S^{-1}P$; para que $P$ sea $A$-plano, es necesario y suficiente
que sea $S^{-1}A$-plano.
\end{env}

\begin{env}[6.3.2]
\label{0.6.3.2-fr}
Sean $A$ un anillo, $B$ una $A$-álgebra, $T$ una parte multiplicativa de
$B$. Si $P$ es un $B$-módulo que es $A$-\emph{plano}, $T^{-1}P$ es $A$-\emph{plano}. En
efecto, para todo $A$-módulo $N$, se tiene
\[
  (T^{-1}P)\otimes_A N
  =(T^{-1}B\otimes_B P)\otimes_A N
  =T^{-1}B\otimes_B(P\otimes_A N)
  =T^{-1}(P\otimes_A N)~;
\]
ahora bien, $T^{-1}(P\otimes_A N)$ es un funtor exacto en $N$, por ser compuesto de los
dos funtores exactos $P\otimes_A N$ (en $N$) y $T^{-1}Q$ (en $Q$). Si $S$
es una parte multiplicativa de $A$ cuya imagen en $B$ está
\emph{contenida en $T$}, $T^{-1}P$ es igual a $S^{-1}(T^{-1}P)$, por tanto es
también $S^{-1}A$-plano en virtud de (6.3.1).
\end{env}

\begin{env}[6.3.3]
\label{0.6.3.3-fr}
Sean $\varphi:A\to B$ un homomorfismo de anillos, $M$ un $B$-módulo. Las
propiedades siguientes son equivalentes:
\begin{enumerate}
  \item[$a)$] $M$ es un $A$-módulo plano.
  \item[$b)$] Para todo ideal maximal $\mathfrak{n}$ de $B$,
    $M_\mathfrak{n}$ es un $A$-módulo plano.
  \item[$c)$] Para todo ideal maximal $\mathfrak{n}$ de $B$, poniendo
    $\mathfrak{m}=\varphi^{-1}(\mathfrak{n})$, $M_\mathfrak{n}$ es un
    $A_\mathfrak{m}$-módulo plano.
\end{enumerate}

En efecto, como $M_\mathfrak{n}=(M_\mathfrak{n})_\mathfrak{m}$,
la equivalencia de $b)$ y $c)$ resulta de (6.3.1), y el hecho de que $a)$
implica $b)$ es un caso particular de (6.3.2). Queda por ver que $b)$
implica $a)$,
\oldpage[0\textsubscript{I}]{57}
es decir, que para todo homomorfismo inyectivo $u:N'\to N$ de
$A$-módulos, el homomorfismo
$v=1\otimes u:M\otimes_A N'\to M\otimes_A N$ es inyectivo. Ahora bien, $v$ es también
un homomorfismo de $B$-módulos, y se sabe que para que sea inyectivo,
basta que para todo ideal maximal $\mathfrak{n}$ de $B$,
\[
  v_\mathfrak{n}:(M\otimes_A N')_\mathfrak{n}
  \longrightarrow(M\otimes_A N)_\mathfrak{n}
\]
sea inyectivo. Pero como
\[
  (M\otimes_A N)_\mathfrak{n}
  =B_\mathfrak{n}\otimes_B(M\otimes_A N)
  =M_\mathfrak{n}\otimes_A N,
\]
$v_\mathfrak{n}$ no es otro que el homomorfismo
\[
  1\otimes u:M_\mathfrak{n}\otimes_A N'
  \longrightarrow M_\mathfrak{n}\otimes_A N,
\]
que es inyectivo puesto que $M_\mathfrak{n}$ es $A$-plano.

En particular (haciendo $B=A$), para que un $A$-módulo $M$ sea plano, es
necesario y suficiente que $M_\mathfrak{m}$ sea $A_\mathfrak{m}$-plano para todo
ideal maximal $\mathfrak{m}$ de $A$.
\end{env}

\begin{env}[6.3.4]
\label{0.6.3.4-fr}
Sea $M$ un $A$-módulo; si $M$ es plano, y si $f\in A$ no es divisor
de $0$ en $A$, $f$ no anula ningún elemento $\neq0$ de $M$, pues
el homomorfismo $m\to f.m$ se escribe $1\otimes u$, donde $u$ es la
multiplicación $a\to f.a$ en $A$ y $M$ está identificado con
$M\otimes_A A$; si $u$ es inyectiva, lo mismo ocurre, pues, con
$1\otimes u$ puesto que $M$ es plano. En particular, si $A$ es
\emph{íntegro}, $M$ es \emph{sin torsión}.

Recíprocamente, supongamos que $A$ es íntegro, $M$ sin torsión, y supongamos que para
todo ideal maximal $\mathfrak{m}$ de $A$, $A_\mathfrak{m}$ es un
\emph{anillo de valoración discreta}; entonces $M$ es $A$-\emph{plano}. Basta, en
efecto (6.3.3), probar que $M_\mathfrak{m}$ es $A_\mathfrak{m}$-plano, y
se puede, pues, suponer que $A$ es ya un anillo de valoración discreta. Pero
como $M$ es límite inductivo de sus submódulos de tipo finito, y estos
últimos son sin torsión, se puede además limitar al caso en que $M$ es de
tipo finito (6.1.2). La proposición resulta en este caso de que $M$ es un
$A$-módulo libre.

En particular, si $A$ es un anillo \emph{íntegro},
$\varphi:A\to B$ un homomorfismo de anillos que hace de $B$ un $A$-módulo
\emph{plano} y $\neq\{0\}$, $\varphi$ es necesariamente \emph{inyectivo}.
Recíprocamente, si $B$ es íntegro, $A$ un subanillo de $B$ y si para todo
ideal maximal $\mathfrak{m}$ de $A$, $A_\mathfrak{m}$ es un anillo de
valoración discreta, $B$ es $A$-\emph{plano}.
\end{env}

